% Class-compatibility preamble retained for a reader excerpt of
%   Sahar Diskin, Anna Geisler, "Perfect Matching in Product Graphs and in their
%   Random Subgraphs", arXiv:2404.14020v4, CC BY 4.0.
%
% The source pins its own publisher class (scrartcl); the delivery wrapper accepts
% only the classes article and amsart, so this file supplies the amsart baseline and
% then copies the author's own declaration block mechanically and unchanged
% (pinned-original-source.tex, source lines 30-57: the author's newtheorem
% declarations and the author's macro declarations).  The author's own package
% list of source lines 4-21 is reused only for the entries the retained text needs
% (enumitem, amsmath, amssymb, amsthm, mathtools, xcolor, geometry) and the
% cross-reference packages already needed by the author's own \Cref/\eqref usage
% (hyperref, cleveref); tikz, caption, subcaption, float, graphicx, wrapfig and
% abstract serve figures and the title block that this excerpt does not carry.
% No mathematical macro, statement, symbol or assumption is introduced, changed
% or removed by this file.
\documentclass[10pt,a4paper]{amsart}
\usepackage[shortlabels]{enumitem}
\usepackage{amsmath}
\usepackage{amssymb}
\usepackage{amsthm}
\usepackage{mathtools}
\usepackage{xcolor}
\usepackage[margin=1in]{geometry}
\usepackage{hyperref}
\usepackage{cleveref}
\hypersetup{
  colorlinks=true,
  linkcolor=blue,
  filecolor=magenta,
  urlcolor=cyan
}

\newtheorem{theorem}{Theorem}
\newtheorem{thm}{Theorem}[section]
\newtheorem{corollary}[thm]{Corollary}
\newtheorem{lemma}[thm]{Lemma}
\newtheorem{proposition}[thm]{Proposition}
\newtheorem{claim}[thm]{Claim}
\newtheorem{remark}[thm]{Remark}
\newtheorem{question}[thm]{Question}

\newcommand{\whp}{whp}
\newcommand{\defi}{\operatorname{def}}
\newcommand{\odd}{\operatorname{odd}}
\newcommand{\iso}{\operatorname{iso}}
\newcommand{\cN}{\mathcal{N}}
\newcommand{\cO}{\mathcal{O}}
\newcommand{\cE}{\mathcal{E}}
\newcommand{\cK}{\mathcal{K}}
\newcommand{\cP}{\mathcal{P}}
\newcommand{\cQ}{\mathcal{Q}}

\newcommand{\footremember}[2]{%
  \footnote{#2}%
  \newcounter{#1}%
  \setcounter{#1}{\value{footnote}}%
}

\newcommand{\ag}[1]{\textcolor{purple}{(Anna: #1)}}
\newcommand{\sd}[1]{\textcolor{red}{(Sahar: #1)}}
\crefname{originalequation}{original equation}{original equations}
\Crefname{originalequation}{Original equation}{Original equations}
\makeatletter
\newlabel{eq:expansion-log}{{1.1}{source}{}{sourcepointer-1}{}}
\newlabel{eq:expansion-log@cref}{{[originalequation][0][]1.1}{[1][0][]source}{}{}{}}
\makeatother
\crefname{originallemma}{original lemma}{original lemmas}
\Crefname{originallemma}{Original lemma}{Original lemmas}
\makeatletter
\newlabel{lem:small-boundary-sets}{{6.3}{source}{}{sourcepointer-2}{}}
\newlabel{lem:small-boundary-sets@cref}{{[originallemma][0][]6.3}{[1][0][]source}{}{}{}}
\makeatother
\title{With high probability there is no small separator whose complement has between u\_1 and n/2 vertices}
\author{Sahar Diskin, Anna Geisler}
\date{arxiv:2404.14020v4; CC BY 4.0}
\begin{document}
\maketitle
\noindent\textit{Editorial note.} All delivered mathematical text is the author's own, copied from the pinned author source sources/combinatorics-fl04-vspace-1-5cm-perfect-matching-in-p/pinned-original-source.tex; no statement or proof step was rewritten, reordered or removed.\par\medskip
\noindent\textit{Editorial note.} Preamble provenance: the source pins the publisher class scrartcl, which the delivery wrapper does not accept, so this card uses sources/combinatorics-fl04-vspace-1-5cm-perfect-matching-in-p/retained-preamble.tex, which supplies the amsart baseline and then carries the author's own newtheorem and macro declarations (pinned source lines 30-57) verbatim. No macro was invented and no mathematical symbol was changed.\par\medskip
\noindent\textit{Editorial note.} Bibliography provenance: the source carries no thebibliography environment and ends with bibliographystyle abbrv plus bibliography perc (author source lines 1691-1692), so the bibliography is resolved with BibTeX from the identical copy staged with the same arXiv version under sources/product-graph-matchings-2404.14020v4/perc.bib; no entry was edited and the author's own abbrv style is used. The counted proof itself contains no citation.\par\medskip
\noindent\textit{Editorial note.} Selection provenance: this card was opened under the file-based control contract (CONTROL.md), with the mechanical pre-card gate run on its statement and proof interval before the card was written and reporting the interval free.\par\medskip
\noindent\textit{Editorial note.} External references are navigation pointers only; every deliberately kept step of the author's argument stays inside the counted spans.\par\medskip
\noindent\textit{Allowed standard prerequisites (editorial).}\begin{itemize}
\item The thresholds u\_0 and u\_1 of the section and their relation, together with the expansion estimate for sets in the product graph. Set up in the author's section on Tutte obstructions and the size thresholds (source lines 655-700); the expansion estimate is the author's displayed inequality eq:expansion-log. Supplies the two size bounds that the proof converts into exponential estimates, and the expansion used in the small-boundary case.
\end{itemize}\medskip
\ifdefined\hypertarget\hypertarget{sourcepointer-1}{}\fi
\noindent\textit{Source reference.} Original equation 1.1: sources/combinatorics-fl04-vspace-1-5cm-perfect-matching-in-p/pinned-original-source.tex, source lines 224--224. Navigation pointer only: the counted proof uses this displayed expansion estimate in its small-boundary case; the display belongs to an earlier section and is not reproduced in this card.\par\medskip
\ifdefined\hypertarget\hypertarget{sourcepointer-2}{}\fi
\noindent\textit{Source reference.} Original lemma 6.3: sources/combinatorics-fl04-vspace-1-5cm-perfect-matching-in-p/pinned-original-source.tex, source lines 714--714. Navigation pointer only: the counted proof uses the author's counting lemma for sets with a small boundary; that lemma has its own card in this lane.\par\medskip
\begin{lemma}\label{lem:no-small-separator}
Whp, there are no disjoint sets $U, B\subseteq V$ satisfying
\[
 |U|\le u_0, \qquad u_1\le|B|\le n/2,
\]
for which every edge from $B$ to $V\setminus(B\cup U)$ is not present in $G_{p_-}$.
\end{lemma}

\begin{proof}
Denote by $b \coloneqq |B|$. The number of possible $U$ is at most
\begin{align*}
    \sum_{u\le u_0}\binom nu \le\exp\left\{u_0\ln(en/u_0)\right\}=\exp\left\{u_0\cdot \frac{M^5}{p_-}\cdot \ln d \right\}\le \exp\left\{\frac{b}{\ln d}\right\},
\end{align*}
where we used the definition of $u_0$ and $u_1$ and that $b\ge u_1$.

If $e_G(B,B^c)\ge b\ln^2 d$, then at least $b\ln^2d-d|U|\ge b\ln^2 d/2$ edges from $B$ avoid $U$. Thus, the probability that there are sets $ U, B$ with $b=|B|$ as in the claim is at most
\[
 \binom nb\exp\left\{\frac{b}{\ln d}\right\}(1-p_-)^{b\ln^2 d/2}=o(n^{-3}).
\]
If $e_G(B,B^c)<b\ln^2 d$, \Cref{lem:small-boundary-sets} gives at most $\exp\{2b/\ln d\}$ choices for $B$.
Moreover the expansion property in \eqref{eq:expansion-log} gives $e_G(B,B^c)\ge \alpha_M b$ for some constant $\alpha_M$ while $du_0=o(b)$, thus at least $\frac{\alpha_M}{2} b$ edges from $B$ avoid $U$. Thus, the probability that there are sets $ U, B$ with $b=|B|$ as in the claim is at most
\[
 \exp\left\{\frac{3b}{\ln d}\right\}(1-p_-)^{\alpha_M b/2}\le\exp\{-\alpha'_M b\}.
\]
A union bound over $b \geq u_1$ concludes the proof.
\end{proof}

\bibliographystyle{plain}
\bibliography{perc}
\end{document}
